\newpage
\scp{\tsc{\ili{Bima}-Lembata}}{B-01}

\noindent{\wg\st{2pt}\begin{xltabular}{\linewidth}
{lllll>{\raggedright\arraybackslash}X}
\lsptoprule
%Language	&	ISO	&	Glotto.	&	banana	&	Etymon	&	Source (LaTeX)	\\	\midrule
%\endfirsthead
Language	&	ISO	&	Glotto.	&	banana	&	Etymon	&	Source	\\	\midrule\hhline{~}
\endhead
\ili{Sanggar}	&	{-}{-}{-}	&	sang1334	&	punti	&	*punti	&	\cite[19]{Munir2017}, \cite[199]{Zollinger1850}	\\
\ili{Bima}	&	bhp	&	bima1247	&	kalo	&	*kalauR	&	DD	\\
\ili{Komodo}	&	kvh	&	komo1261	&	kalo	&	*kalauR	&	DD	\\
W. \ili{Manggarai}	&	(mqy)	&	west2537	&	muku	&	{\#}muku	&	\cite[338]{Verheijen1967}	\\
WC. \ili{Manggarai}	&	(mqy)	&	west2536	&	muku/ 	&	{\#}muku/	&	\cite[338, 519]{Verheijen1967}	\\	\hhline{~}
&&&punti\footnotemark[7]&*punti&\hp{\,}\\
C. \ili{Manggarai}	&	(mqy)	&	cent2073	&	muku	&	{\#}muku	&	\cite[338]{Verheijen1967}	\\
E. \ili{Manggarai}	&	(mqy)	&	east2463	&	muku	&	{\#}muku	&	\cite[338]{Verheijen1967}	\\
\ili{Riung}	&	riu	&	riun1237	&	muku	&	{\#}muku	&	\cite[338]{Verheijen1967}	\\
\ili{Rembong}	&	reb	&	remb1250	&	muku	&	{\#}muku	&	\cite{Verheijen1977}	\\
\ili{Waerana}	&	wrx	&	waer1237	&	muku	&	{\#}muku	&	\cite{VerheijenDadu1994}	\\
Palu{\Q}e	&	ple	&	palu1252	&	muku	&	{\#}muku	&	\cite{Donohue2003}	\\
\ili{Lio}	&	ljl	&	lioo1240	&	muku	&	{\#}muku	&	\cite{Elias2017}	\\
\ili{Ende}	&	end	&	ende1246	&	muku	&	{\#}muku	&	DD	\\
\ili{Rongga}	&	ror	&	rong1269	&	muku	&	{\#}muku	&	DD	\\
So{\Q}a	&	ssq	&	soaa1237 	&	muku	&	{\#}muku	&	\cite{Elias2017}	\\
\ili{Ngadha}	&	nxg	&	ngad1261	&	muku	&	{\#}muku	&	DD	\\
E. \ili{Nage}	&	nxe	&	east2838	&	muku	&	{\#}muku	&	\cite{Elias2017}	\\
Nga{\Q}o, \ili{Oja}	&	(nxe)	&	ngao1245	&	muku	&	{\#}muku	&	\cite{Elias2017}	\\
\ili{Keo}	&	xxk	&	keoo1238	&	muku	&	{\#}muku	&	\cite{Baird2002}	\\
\ili{Sika}	&	ski	&	sika1262	&	muʔu	&	{\#}muku	&	DD	\\
\ili{Kedang}	&	ksx	&	keda1252	&	muʔu	&	{\#}muku	&	DD	\\
\ili{Ile Mandiri} &	(slp)	&	nucl1534	&	muko	&	{\#}muku	&	\cite{Klamer2015}	\\	\hhline{~}
\ili{Lamaholot}&&&&&\hp{\,}\\
\ili{Solor}	&	(slp)	&	west2543	&	muko	&	{\#}muku	&	\cite[326]{Kroon2016}	\\	\hhline{~}
\ili{Lamaholot}&&&&&\hp{\,}\\
Adonara	&	adr	&	adon1237	&	muko	&	{\#}muku	&	\cite{Klamer2015}	\\	\hhline{~}
\ili{Lamaholot}&&&&&\hp{\,}\\
\ili{Alorese}	&	aol	&	alor1247	&	muko	&	{\#}muku	&	\cite{Sulistyono2018}	\\
\ili{Kalikasa}	&	lvu	&	levu1239	&	muku,	&	{\#}muku	&	\cite{Fricke2015b}	\\	\hhline{~}
&&&mukor&&\hp{\,}\\
\end{xltabular}}
\footnotetext[7]{\citet[519]{Verheijen1967} gives \ili{Pacar} \ili{Manggarai} \it{punti} ‘k.o.
			sweet banana with small,
			short fruit’.
			%(‘sb pisang; buahnya kecil pendek, manis’).
			\citet[80]{Verheijen1984} states regarding this term: “I connected
			it with Dempwolff’s AN *pun(t)i(h),
			By now I think that this name without cognates in the very neighbourhood
			and locally so limited must be regarded as a loan.
			Possibly a Menadonese official was the introducer.”}

\noindent{\wg\st{2pt}\begin{xltabular}{\linewidth}
{lllll>{\raggedright\arraybackslash}X}
\lsptoprule
Language	&	ISO	&	Glotto.	&	banana	&	Etymon	&	Source	\\	\midrule\hhline{~}
\endhead
\ili{Painara}	&	lmf	&	sout2896	&	mukor	&	{\#}muku	&	\cite{Fricke2015a}	\\
%&&&&&\\\hhline{~}
\ili{Kodi}	&	kod	&	kodi1247	&	kaloɣo/ 	&	*kalauR/	&	\cite[149, 293]{Onvlee1984}	\\	\hhline{~}
&&&muku\footnotemark[8]&{\#}muku&\hp{\,}\\
\ili{Laboya}	&	lmy	&	lamb1273	&	kaloːwo	&	*kalauR	&	\cite{Verdizade2019b}	\\
\ili{Weyewa}	&	wew	&	weye1237	&	kalowu/	&	*kalauR	&	\cite[149, 293]{Onvlee1984}	\\	\hhline{~}
&&&kalowo&&\hp{\,}\\	\hhline{~}
&&&kaka&&\hp{\,}\\
S. \ili{Weyewa}	&	(wew)	&		&	‹ka.lo.wo›	&	*kalauR	&	\cite[vol. 2, 307]{Putra2007}	\\
\ili{Tana Righu}	&	(wew)	&	tana1282	&	‹ka.lo.ɣo›	&	*kalauR	&	\cite[vol. 2, 307]{Putra2007}	\\
\ili{Loli}	&	(wew)	&	laul1237	&	kalowo	&	*kalauR	&	\cite{Mitchell2008a}	\\
\ili{Loura}	&	lur	&	nucl1519	&	‹ka.lo.ɣo›	&	*kalauR	&	\cite[vol. 2, 307]{Putra2007}	\\
\ili{Mamboru}	&	mvd	&	mamb1305	&	kalowu/&	*kalauR/	&	\cite[149, 293]{Onvlee1984}	\\	\hhline{~}
&&&muku&{\#}muku&\hp{\,}\\
\ili{Anakalang}	&	akg	&	anak1240	&	kalowu/	&	*kalauR/	&	\cite[149, 293]{Onvlee1984}	\\	\hhline{~}
&&&muku&{\#}muku&\hp{\,}\\
\ili{Wanukaka}	&	wnk	&	wanu1241	&	kaloʔa	&	*kalauR	&	DD	\\
\ili{Mangili}	&	(xbr)	&	mang1406	&	kaluu/	&	*kalauR/	&	\cite[149, 293]{Onvlee1984}	\\	\hhline{~}
&&&muku&{\#}muku&\hp{\,}\\
\ili{Kambera}	&	(xbr)	&	kamb1299	&	kaluu/	&	*kalauR/	&	\cite[149, 293]{Onvlee1984}	\\	\hhline{~}
&&&muku&{\#}muku&\hp{\,}\\
\ili{Melolo}/\ili{Rindi}	&	(xbr)	&	melo1243	&	‹kà.lu›	&	*kalauR	&	\cite[vol. 2, 307]{Putra2007}	\\
Tabundung-	&	(xbr)	&		&	‹kà.lu›	&	*kalauR	&	\cite[vol. 2, 307]{Putra2007}	\\	\hhline{~}
Karera&&&&&\hp{\,}\\
\ili{Lewa}	&	(xbr)	&	lewa1240	&	kalowu/	&	*kalauR/	&	\cite[149, 293]{Onvlee1984}	\\	\hhline{~}
&&&muku&{\#}muku&\hp{\,}\\
Kanatang	&	(xbr)	&	kana1287	&	‹kà.lu›	&	*kalauR	&	\cite[vol. 2, 307]{Putra2007}	\\
\ili{Prai Paha}	&	(xbr)	&		&	‹ka.lo›	&	*kalauR	&	\cite[vol. 2, 307]{Putra2007}	\\
Umbu Ratu 	&	(xbr)	&	umar1237	&	‹ka.lo.wu›	&	*kalauR	&	\cite[vol. 2, 307]{Putra2007}	\\	\hhline{~}
Nggai&&&&&\hp{\,}\\
\ili{Dhao}	&	nfa	&	dhao1237	&	muʔu	&	{\#}muku	&	Grimes fieldnotes	\\
\ili{Havu}	&	hvn	&	sabu1255	&	muʔu	&	{\#}muku	&	Grimes fieldnotes	\\
\lspbottomrule
\end{xltabular}}
\footnotetext[8]{\citet[149]{Onvlee1984} gives \ili{Kambera} \it{kaluu} (with \ili{Mangili},
			\ili{Lewa}, \ili{Anakalang}, \ili{Mamboru}, \ili{Weyewa},
			and \ili{Kodi} cognates) as a generic term for ‘banana’.
			\ili{Kambera} \it{muku}, (as well as cognates
			and \ili{Weyewa} \it{kalowo kaka}) is ‘kind of banana which is good
			to fry’.}\setcounter{footnote}{8}